\chapter{Angioplasty}
\index{Angioplasty}

\section*{Definition and Overview}
Angioplasty is a minimally invasive procedure that reopens a narrowed or blocked blood vessel. A small balloon is inflated inside the vessel at the site of the narrowing, crushing the underlying plaque against the vessel wall and restoring blood flow; a mesh tube called a stent is usually left in place to keep the vessel open. The procedure is performed through a small puncture in the groin or wrist, so it avoids open surgery.

The best-known form is coronary angioplasty, which treats the arteries of the heart and is called a percutaneous coronary intervention (PCI). The same technique is applied to arteries of the legs, kidneys, and neck, and to bypass grafts and dialysis fistulas. Angioplasty is used in two main situations: as emergency treatment during a heart attack, where rapid reopening of the blocked artery saves heart muscle, and as elective treatment to relieve angina or claudication when symptoms are not controlled by medicines.

\section*{Epidemiology}
Coronary angioplasty is one of the most commonly performed procedures worldwide, with millions of PCI procedures carried out each year, predominantly for acute heart attacks and for significant angina. Peripheral angioplasty of the leg arteries is increasingly common in older adults, especially in smokers and people with diabetes, among whom peripheral arterial disease is frequent.

Rates of the procedure have risen steadily as techniques and stent technology have improved and as cardiac catheter laboratories have become more widely available. At the same time, the balance between angioplasty, bypass surgery, and medical therapy is now tailored to the individual, with modern trials showing that many people with stable angina do just as well with medicines as with intervention.

\section*{Aetiology and Pathophysiology}
Angioplasty treats the consequences of vessel narrowing and blockage, most commonly caused by:

\begin{itemize}
    \item \textbf{Atherosclerosis:} fatty and fibrous plaques accumulate in the arterial wall, narrowing the lumen and restricting blood flow
    \item \textbf{Plaque rupture and thrombosis:} in a heart attack, a vulnerable plaque breaks open and a blood clot forms on top, suddenly occluding the artery
    \item \textbf{Coronary heart disease:} narrowing of the heart arteries that causes stable or unstable angina and heart attack
    \item \textbf{Peripheral arterial disease:} narrowing of the leg arteries causing cramping leg pain on walking (intermittent claudication) or critical limb ischaemia
    \item \textbf{Renovascular disease:} renal artery narrowing that can contribute to difficult-to-control hypertension and declining kidney function
    \item \textbf{Carotid artery disease:} narrowing of the neck arteries, an important risk factor for stroke
    \item \textbf{Failed grafts and fistulas:} angioplasty can reopen blocked coronary bypass grafts and failing dialysis fistulas
\end{itemize}

The mechanical principle is straightforward: inflating a balloon fractures and remodels the plaque, enlarges the lumen, and stretches the vessel wall; the stent then provides a scaffold that prevents elastic recoil and seals dissection flaps, maintaining the improved lumen after the balloon is deflated.

\section*{Clinical Features}
Angioplasty itself causes few symptoms, and the experience reflects the procedure and the condition being treated:

\begin{itemize}
    \item Performed under local anaesthetic at the puncture site, usually with light sedation
    \item A catheter and guidewire are threaded through the artery to the blockage under X-ray guidance
    \item A small balloon is inflated at the narrowing for a few seconds at a time, crushing the plaque and stretching the vessel
    \item A stent is expanded to hold the artery open and is left in place permanently
    \item Contrast dye is used throughout, producing a warm, flushing sensation when injected
    \item The procedure typically takes 30--90 minutes, depending on its complexity
    \item Symptoms such as angina or leg cramping usually improve promptly once blood flow is restored
\end{itemize}

\section*{Differential Diagnosis}
Angioplasty is chosen over alternative treatments on the basis of the angiogram, symptoms, and overall health:

\begin{itemize}
    \item Significant narrowing treatable with angioplasty versus minor narrowing managed with medicines alone
    \item Angioplasty versus coronary bypass surgery (CABG), a decision based on the number, site, and complexity of the narrowed arteries
    \item Angioplasty versus medical therapy in stable angina, which is a reasonable first choice for many people
    \item A localised, treatable blockage versus diffuse disease too widespread for the procedure
    \item Leg artery narrowing versus other causes of leg pain, such as nerve root or joint disease
    \item Carotid stenosis requiring intervention versus smaller lesions best treated with medicines alone
\end{itemize}

\section*{When to Seek Medical Care}
After angioplasty, contact the hospital or your doctor if you develop:

\begin{itemize}
    \item Heavy or growing bleeding, bruising, or a painful swelling at the puncture site
    \item Coldness, numbness, tingling, or a colour change in the treated limb
    \item Chest pain or breathlessness that does not settle
    \item Fever, chills, or spreading redness at the puncture site
    \item Worsening of the symptoms of the original condition
\end{itemize}

\textbf{Contact your local emergency services for:}
\begin{itemize}
    \item Sudden severe chest pain, especially with sweating, nausea, or breathlessness
    \item Sudden weakness, numbness, or difficulty speaking
    \item Heavy bleeding from the puncture site that will not stop with pressure
    \item Fainting or collapse
\end{itemize}

\section*{Diagnosis and Investigations}
Preparation for angioplasty mirrors that for an angiogram and includes:

\begin{itemize}
    \item \textbf{Coronary or peripheral angiography:} the images show exactly where, and how severely, the vessels are narrowed
    \item \textbf{Blood tests:} kidney function, clotting screen, full blood count, and cardiac enzymes (troponin) in acute settings
    \item \textbf{ECG:} to assess rhythm, ischaemia, or the changes of a heart attack
    \item \textbf{Echocardiography:} to assess heart function in selected cases
    \item \textbf{Medication review:} especially blood thinners, diabetes medicines, and kidney function, which affects contrast dye safety
    \item \textbf{Informed consent:} a discussion of the benefits, risks, and alternatives, including medical therapy and surgery
\end{itemize}

Additional physiological measurements, such as the pressure gradient across a narrowing (fractional flow reserve), are sometimes used during the procedure to confirm that a stenosis is significant enough to justify stenting.

\section*{Management}
\begin{itemize}
    \item \textbf{Antiplatelet therapy:} aspirin plus a second agent such as clopidogrel is taken before and for a defined period after stenting to prevent stent thrombosis
    \item \textbf{Balloon and stent:} the narrowed segment is dilated and scaffolded
    \item \textbf{Drug-eluting stents:} stents coated with a slow-releasing medicine that reduces re-narrowing (in-stent restenosis)
    \item \textbf{Specialised balloons:} drug-coated or cutting balloons for recurrent or calcified narrowings
    \item \textbf{Clot extraction:} in a heart attack, thrombus may be aspirated before ballooning
    \item \textbf{Aftercare:} lie flat for several hours, keep the puncture limb still, and drink fluids to protect the kidneys
    \item \textbf{Cardiac rehabilitation:} a supervised exercise and education programme after coronary angioplasty
    \item \textbf{Long-term medicines:} statins and blood pressure or anti-anginal drugs as required, regardless of the procedure
\end{itemize}

\section*{Complications}
\begin{itemize}
    \item Bleeding, bruising, or a pseudoaneurysm (false aneurysm) at the puncture site
    \item Damage to the treated artery, including dissection (a tear) or spasm
    \item Contrast dye allergy and contrast-induced kidney injury
    \item Stent thrombosis -- a blood clot forming on the stent that can occlude the vessel
    \item In-stent restenosis -- re-narrowing inside the stent over months
    \item Injury to the heart muscle, abnormal heart rhythms, or stroke (all rare)
    \item Heart attack or the need for emergency bypass surgery in a small proportion of cases
\end{itemize}

\section*{Prognosis and Outcomes}
Angioplasty is highly effective at restoring blood flow and relieving symptoms. In a heart attack, prompt angioplasty of the blocked artery reduces heart muscle damage and improves survival more effectively than clot-busting drugs alone. In stable angina, it reliably relieves chest pain and improves exercise capacity, although trials show it does not clearly reduce the long-term risk of death compared with optimal medical therapy. Re-narrowing is uncommon with modern drug-eluting stents, and most people return to normal activities within a week. As with all vascular disease, lifestyle changes and lifelong medication remain essential to the long-term outlook.

\section*{Prevention and Self-Care}
\begin{itemize}
    \item Do not smoke -- stopping is the single most effective step
    \item Keep blood pressure, cholesterol, and blood sugar well controlled
    \item Eat a heart-healthy diet and maintain a healthy weight
    \item Be physically active within a plan agreed with your treating team
    \item Take antiplatelet and other medicines exactly as prescribed, and never stop blood thinners without medical advice
    \item Attend cardiac rehabilitation and scheduled follow-up appointments
    \item Know the warning signs of a heart attack and seek emergency help early
\end{itemize}

\section*{Sources and Further Reading}
The following sources provide reliable further information on angioplasty:

\begin{itemize}
    \item NHS. \textit{Information on coronary angioplasty and stent insertion.} https://www.nhs.uk/conditions/coronary-angioplasty/
    \item British Heart Foundation. \textit{Plain-language guide to angioplasty, stents, and recovery.} https://www.bhf.org.uk/informationsupport
    \item MedlinePlus. \textit{Overview of angioplasty and stent placement, including preparation and recovery.} https://medlineplus.gov/angioplasty.html
    \item Mayo Clinic. \textit{Coronary angioplasty and stents: purpose, risks, and recovery.} https://www.mayoclinic.org/tests-procedures/angioplasty/about/pac-20384761
    \item MSD Manual. \textit{Percutaneous coronary intervention in the professional health section.} https://www.msdmanuals.com/professional/cardiovascular-disorders/coronary-artery-disease/treatment-of-coronary-artery-disease
\end{itemize}
